%
% 4-rothsteintrager.tex
%
% (c) 2026 Prof Dr Andreas Müller
%
\section{Die Methode von Rothstein-Trager
\label{buch:ratint:section:rothsteintrager}}
\kopfrechts{Die Methode von Rothstein-Trager}%
Die Methode von Rothstein-Trager ermöglicht, auch den logarithmischen
Teil
\begin{equation}
I
=
\int
\frac{p(z)}{q(z)}
\,dz
\label{buch:ratint:rothsteintrager:eqn:integral}
\end{equation}
für einen quadratfreien Nenner $q(z)$ und einen dazu teilerfremden
Zähler $p(z)$ mit $\deg p(z)<\deg q(z)$ zu berechnen.

%
% Residuen und die Form des Integrals
%
\subsection{Residuen und die Form des Integrals
\label{buch:ratint:rothsteintrager:subsection:residuen}}
In den komplexen zahlen lässt sich der quadratfreie Nenner $q(z)$ von
\eqref{buch:ratint:rothsteintrager:eqn:integral}
in ein Produkt verschiedner Linearfaktoren
\begin{equation}
q(z)
=
(z-\alpha_1)\cdots(z-\alpha_n)
\label{buch:ratint:rothsteintrager:eqn:faktorisierung}
\end{equation}
zerlegen.
Da die Nullstellen alle verschieden sind, hat die Partialbruchzerlegung
des Integranden die Form
\begin{equation}
\frac{p(z)}{q(z)}
=
\frac{A_1}{z-\alpha_1}
+
\ldots
+
\frac{A_n}{z-\alpha_n}.
\label{buch:ratint:rothsteintrager:eqn:integrand}
\end{equation}
In dieser Form lässt sich das Integral jedoch sofort angeben, es ist
\begin{equation}
\int
\frac{p(z)}{q(z)}
\,dz
=
A_1\log(z-\alpha_1)
+
\ldots
+
A_n\log(z-\alpha_n).
\label{buch:ratint:rothsteintrager:eqn:integralform}
\end{equation}
Der Ausdruck kann noch etwas vereinfacht werden, wenn zwei der
Koeffizienten gleich sind.
Ist nämlich $A_k=A_l$, dann können die beiden Summanden in
$A_k\log\bigl((z-\alpha_k)(z-\alpha_k)\bigr)$
zusammengefasst werden.
Für entgegengesetzte Koeffizienten $A_k=-A_l$ folgt entsprechend
\[
A_k\log \frac{z-\alpha_k}{z-\alpha_l}.
\]
Diese Form tritt zum Beispiel bei der Integration reeller Funktionen
auf, wo die Nullstellen des Nenners paarweise konjugiert komplex sind.

Wenn die Faktorisierung
\eqref{buch:ratint:rothsteintrager:eqn:faktorisierung}
von $q(z)$ bekannt ist, ist das Integrationsproblem an dieser
Stelle gelöst.
Wie früher bereits dargelegt wurde, kann man jedoch nicht davon
ausgehen, dass das Faktorisierungsproblem gelöst werden kann.
Es müssen also alternative Methoden gesucht werden, die ohne die 
die Faktorisierung auskommen.

%
% Die Koeffizienten A_k und Residuen
%
\subsubsection{Die Koeffizienten $A_k$ und Residuen}
Die Form 
\eqref{buch:ratint:rothsteintrager:eqn:integrand}
des Integranden des logarithmischen Teils zeigt, dass er isolierte
Singularitäten hat und daher um jede Singularität in eine Laurent-Reihe
entwickelt werden kann.
Man kann sogar sagen, dass die Singularitäten Pole erster Ordnung
sein müssen.

Sie jetzt $\alpha$ eine dieser Pole, dann kann der Integrand
in eine Laurent-Reihe
\[
\frac{p(z)}{q(z)}
=
\frac{b_1}{z-\alpha}
+
a_0+a_1(z-\alpha) + a_2(z-\alpha)^2+\ldots
\]
um den Punkt $\alpha$ entwickelt werden.
Gliedweise Integration der Reiher ergibt in einer Umgebung von $\alpha$
\[
\int
\frac{p(z)}{q(z)}
\,dz
=
b_1\log(z-\alpha)
+
a_0(z-\alpha) + \frac{a_1}{2}(z-\alpha)^2 + \frac{a_2}{3}(z-\alpha)^3
+\ldots
\]
Durch Vergleich mit
\eqref{buch:ratint:rothsteintrager:eqn:integralform}
folgt, dass die gesuchten Koeffizienten $A_k$ die Residuen
des Integranden an den Nullstellen von $q(z)$ sein müssen.
Diese Koeffizienten lassen sich aber direkt aus den Polynomen 
$p(z)$ und $q(z)$ berechnen.

\begin{satz}
\label{buch:ratint:rothsteintrager:satz:residuum}
Sind $p(z)$ und $q(z)$ teilerfremd und $q(z)$ quadratfrei und
$\alpha$ eine Nullstelle des Nenners.
Dann ist das Residuum von $p(z)/q(z)$ an der Stelle $\alpha$
\[
A
=
\operatorname{res}\biggl(\alpha;\frac{p(z)}{q(z)}\biggr)
=
\frac{p(\alpha)}{q'(\alpha)}.
\]
\end{satz}

Der Satz ist eine Folge des folgenden, algemeineren Satzes.
Er gilt für einen beliebigen Bruch von holomorphen Funktionen,
dessen Nenner eine einfache Nullstelle hat.

\begin{satz}
Ist $f(z)$ in einer Umgebung von $\alpha$ holomorph und hat $g(z)$
eine Nullstelle erster Ordnung in $\alpha$, dann ist das Residuum
von $f(z)/g(z)$ an der Stelle $\alpha$ der Wert $f(\alpha)/g'(\alpha)$.
\end{satz}

\begin{proof}
Da $g(z)$ eine Nullstelle erster Ordnung an der Stelle $\alpha$
hat, kann man
\[
g(z) = (z-\alpha)\bigl(g'(\alpha)+ r(z)(z-\alpha)\bigr)
\]
und damit
\[
\frac{f(z)}{g(z)}
=
\frac{f(z)}{(z-\alpha)\bigl(g'(\alpha)+r(z)(z-\alpha)\bigr)}
=
\frac{1}{z-\alpha}
\cdot
\frac{f(z)}{g'(\alpha)+r(z)(z-\alpha)}
=
\frac{1}{z-\alpha}
\cdot
h(z)
\]
mit einer in einer Umgebung von $\alpha$ holomorphen Funktion
$h(z)$ schreiben.
Nach Satz~\ref{buch:analytisch:laurent:satz:laurentreihe} ist
das Residuum das Integral
\begin{align*}
\frac{1}{2\pi i}
\oint_\gamma
\frac{f(z)}{g(z)}
\,dz
&=
\frac{1}{2\pi i}
\oint_\gamma
\frac{1}{z-\alpha}
\frac{f(z)}{g'(\alpha)+r(z)(z-\alpha)}
\,dz
\\
&=
\frac{1}{2\pi i}
\oint_\gamma
\frac1{z-\alpha}
\cdot
h(z)
\,dz.
\intertext{Nach dem Integralsatz~\ref{buch:integration:satz:residuensatz}
ergibt das Integral}
&=
h(\alpha)
=
\frac{f(\alpha)}{g'(\alpha)}.
\end{align*}
Dies beweist den Satz.
\end{proof}

%
% Der Satz von Rothstein-Trager
%
\subsection{Der Satz von Rothstein-Trager}
Die Form des logarithmischen Teils des Integrals ist zwar bekannt,
und der Satz~\ref{buch:ratint:rothsteintrager:satz:residuum}
ermöglicht eine direkte Berechnung des Koeffizienten, sobald
die Nullstellen bekannt sind, aber genau die Bestimmung der
Nullstellen war ja der Punkt, warum die Methode von Bernoulli-Methode
als unzureichend erkannt wurde.
Dies ändert der folgende Satz von Rothstein-Trager.

\begin{satz}[Rothstein-Trager]
\label{buch:ratint:rothsteintrager:satz:rothsteintrager}
Ist $p(z)$ teilerfremd zum quadratfreien Polynom $q(z)$ mit
$\deg p(z)<\deg q(z)$, dann ist
\[
\int
\frac{p(z)}{q(z)}
\,dz
=
\sum_{i=1}^m c_i \log v_i(z),
\]
wobei die $c_i$ die möglicherweise komplexen Zahlen $c$ sind, derart,
dass $p(z)-cq'(z)$ und $q(z)$ einen gemeinsamen Faktor haben.
Die zugehörigen $r_i(z)$ sind
\[
r_i(z)
=
\operatorname{ggT}\bigl(p(z)-c_iq'(z),q(z)\bigr).
\]
\end{satz}

\begin{proof}
TODO
\end{proof}

Wenn es gelingt, eine algebraische Methode zur Bestimmung der
Koeffizienten $c_i$ zu finden, können auch die Faktoren 
$r_i(z)$ gefunden werden.
Man beachte, dass nicht mehr eine Faktorisierung von $q(z)$
verlangt wird, sondern nur noch die Möglichkeit, gemeinsame
Faktoren zu finden, was im Prinzip mit dem euklidischen Algorithmus
möglich ist.

%
% Gemeinsame Nullstellen und die Resultante
%
\subsubsection{Gemeinsame Nullstellen und die Resultante}
Der Satz~\ref{buch:ratint:rothsteintrager:satz:rothsteintrager}
verlagert das Problem, die Nullstellen des Nenners zu finden darauf,
die Koeffizienten $c_i$ zu finden.
Dies ist aus zwei Gründen einfacher:
\begin{enumerate}
\item
Oft sind einige der Koeffizienten $A_k$ gleich.
In diesen Fällen ist es nicht nötig, die zugehörigen Linearfaktoren
zu kennen, es genügt das Produkt dieser Linearfaktoren zu finden,
dies ist das Polynom $r_i(z)$ im
Satz~\ref{buch:ratint:rothsteintrager:satz:rothsteintrager}.
\item
Für das Kriterium, einen gemeinsamen Teiler zu finden gibt es
ein einfaches algebraisches Kriterium, nämlich die Resultante.
\end{enumerate}

In Kapitel~\ref{chapter:resultante} wird beschrieben, wie die 
Resultante des Problem, die Koeffizienten $c_i$ zu finden, lösen
kann.

